\documentclass{article}
\usepackage{amsmath,amssymb,amsthm}
\usepackage{algorithm}
\usepackage{algpseudocode}
\usepackage{hyperref}
\usepackage{enumitem}
\newtheorem{theorem}{Theorem}
\newtheorem{lemma}{Lemma}
\newtheorem{assumption}{Assumption}
\newcommand{\E}{\mathbb{E}}
\newcommand{\R}{\mathbb{R}}

\begin{document}

\section*{AdaGrad with Scalar Accumulator}

\section*{Mathematical Formulation}
Let $f:\R^d\to\R$ be a convex and $L$‑smooth objective.  
Denote by $g_k:=\nabla f(w_k)$ the (full) gradient at iteration $k$.  
The algorithm keeps a scalar (or diagonal) accumulator of squared gradients  

\[
G_k \;:=\; \epsilon+\sum_{i=1}^{k} g_i\odot g_i\in\R_{+}^{d},
\qquad (\epsilon>0\text{ for numerical stability}).
\]

The learning‑rate vector at step $k$ is  

\[
\eta_k \;:=\; \frac{\alpha}{\sqrt{G_k}}\in\R_{+}^{d},
\]

where $\alpha>0$ is a user‑chosen base step size.  
The update rule is therefore  

\[
\boxed{ \; w_{k+1}= w_k - \eta_k\odot g_k
      = w_k - \frac{\alpha}{\sqrt{G_k}}\odot g_k \; } .
\]

For the scalar case $d=1$ this reduces to  

\[
G_k = \epsilon +\sum_{i=1}^{k} g_i^{2},\qquad 
w_{k+1}= w_k - \frac{\alpha}{\sqrt{G_k}}\,g_k .
\]

\section*{Pseudocode}
\begin{algorithm}[H]
\caption{AdaGrad (scalar/diagonal version)}\label{alg:adagrad}
\begin{algorithmic}[1]
\Require Initial point $w_0\in\R^{d}$, base step size $\alpha>0$, 
        stability constant $\epsilon>0$.
\State $G_0 \gets \epsilon\mathbf{1}_d$ \Comment{initialize accumulator}
\For{$k=0,1,2,\dots,T-1$}
    \State $g_k \gets \nabla f(w_k)$
    \State $G_{k+1} \gets G_k + g_k\odot g_k$
    \State $\eta_k \gets \alpha / \sqrt{G_{k+1}}$ \Comment{elementwise division}
    \State $w_{k+1} \gets w_k - \eta_k\odot g_k$
\EndFor
\State \Return $w_T$ (or the average $\bar w_T=\tfrac1T\sum_{k=1}^{T}w_k$)
\end{algorithmic}
\end{algorithm}

\section*{Dry Run Example}
Consider the one‑dimensional quadratic  

\[
f(w) = (w-3)^2 ,\qquad \nabla f(w)=2(w-3).
\]

Parameters: $w_0=0$, $\alpha=0.1$, $\epsilon=0$ (for illustration).  

\begin{center}
\begin{tabular}{c|c|c|c|c}
$k$ & $w_k$ & $g_k=2(w_k-3)$ & $G_k=\sum_{i=1}^{k} g_i^2$ & $w_{k+1}=w_k-\frac{\alpha}{\sqrt{G_k}}g_k$\\\hline
0 & $0$ & $-6$ & $36$ & $0-\frac{0.1}{\sqrt{36}}(-6)=0+0.1\cdot\frac{6}{6}=0.1$\\
1 & $0.1$ & $2(0.1-3)=-5.8$ & $36+33.64=69.64$ &
$0.1-\frac{0.1}{\sqrt{69.64}}(-5.8)
=0.1+0.1\cdot\frac{5.8}{8.34}\approx0.1694$\\
2 & $0.1694$ & $2(0.1694-3)=-5.6612$ &
$69.64+32.06\approx101.70$ &
$0.1694-\frac{0.1}{\sqrt{101.70}}(-5.6612)
\approx0.1694+0.1\cdot\frac{5.6612}{10.08}\approx0.2322$
\end{tabular}
\end{center}

Corresponding objective values $f(w_k)$: $9$, $8.60$, $8.31$, $8.05$, showing a monotone decrease.

\section*{Assumptions}
\begin{assumption}[Convexity and Smoothness]\label{ass:convex}
$f:\R^{d}\to\R$ is convex and differentiable, and $L$‑smooth:
\[
\|\nabla f(x)-\nabla f(y)\|\le L\|x-y\|,\quad \forall x,y\in\R^{d}.
\]
\end{assumption}
\begin{assumption}[Bounded Subgradients]\label{ass:bound}
There exists $G_{\max}>0$ such that $\|g_k\|=\|\nabla f(w_k)\|\le G_{\max}$ for all $k$.
\end{assumption}
\begin{assumption}[Existence of a Minimizer]\label{ass:min}
Let $w^{\ast}\in\arg\min_{w}f(w)$ and denote $D:=\|w_0-w^{\ast}\|$.
\end{assumption}

\section*{Main Convergence Theorem}
\begin{theorem}[AdaGrad Convergence for Convex $L$‑smooth objectives]\label{thm:adagrad}
Under Assumptions \ref{ass:convex}--\ref{ass:min}, for any $T\ge1$ the averaged iterate 
$\bar w_T:=\frac1T\sum_{k=1}^{T} w_k$ satisfies  

\[
\E\!\bigl[f(\bar w_T)-f(w^{\ast})\bigr]
\;\le\;
\frac{D^2}{2\alpha}\,\frac{1}{\sqrt{T}}
\;+\;
\frac{\alpha G_{\max}^{2}}{2}\,\frac{\log\!\bigl(1+T G_{\max}^{2}/\epsilon\bigr)}{T},
\]
and in particular, choosing $\alpha = D/ G_{\max}\sqrt{\log(1+TG_{\max}^{2}/\epsilon)}$ yields  

\[
\E\!\bigl[f(\bar w_T)-f(w^{\ast})\bigr]=
O\!\left(\frac{G_{\max}D\sqrt{\log T}}{\sqrt{T}}\right).
\]
\end{theorem}

\section*{Theoretical Properties}
\begin{enumerate}[label=\arabic*.]
\item \textbf{Adaptive step sizes.}  The denominator $\sqrt{G_k}$ grows with the accumulated squared gradients, automatically shrinking the step in directions that have seen large gradients.
\item \textbf{Stability.}  Because $\eta_k\le \alpha/\sqrt{\epsilon}$ for all $k$, the algorithm never takes steps larger than the initial base step size.
\item \textbf{Hyper‑parameter sensitivity.}  The bound depends on $\alpha$ only through the product $\alpha G_{\max}$ and the term $D^2/(2\alpha)$, implying that the optimal $\alpha$ balances these two contributions.
\item \textbf{Comparison to SGD.}  For constant step size SGD the optimal rate is $O(1/\sqrt{T})$ with a constant that does not adapt to the geometry of the gradients.  AdaGrad improves the constant by a factor $\sqrt{\log T}$ when the gradients are sparse or have large variance across coordinates.
\end{enumerate}

\section*{Discussion}
The theorem shows that AdaGrad attains the optimal $O(1/\sqrt{T})$ convergence for convex smooth problems while automatically adapting to the observed gradient magnitudes.  The logarithmic factor is unavoidable in the worst case (see lower bounds for adaptive methods).  In practice the method is especially beneficial when some coordinates receive far fewer updates than others, as the per‑coordinate accumulation leads to larger steps on rarely‑updated dimensions.

\section*{Preliminaries (Definitions and Lemmas)}
\begin{lemma}[Three‑point identity for $L$‑smooth convex functions]\label{lem:three}
For any $x,y\in\R^{d}$,
\[
f(y)\le f(x)+\langle\nabla f(x),y-x\rangle+\frac{L}{2}\|y-x\|^{2}.
\]
\end{lemma}
\begin{proof}
Direct consequence of $L$‑smoothness; see standard textbooks. ∎
\end{proof}
\begin{lemma}[Bound on the accumulator]\label{lem:acc}
For all $k\ge1$,
\[
\sqrt{G_k}\;\ge\;\sqrt{\epsilon}\;+\;\frac{1}{\sqrt{G_k}}\sum_{i=1}^{k} \|g_i\|^{2}.
\]
\end{lemma}
\begin{proof}
Since $G_k=\epsilon+\sum_{i=1}^{k}\|g_i\|^{2}$ and all terms are non‑negative,
$\sqrt{G_k}\ge \sqrt{\epsilon}$ and $\sqrt{G_k}\ge \sqrt{\sum_{i=1}^{k}\|g_i\|^{2}}$.
Multiplying the second inequality by $\sqrt{G_k}$ yields the claimed bound. ∎
\end{proof}

\section*{Main Proof (Step‑by‑Step Derivation)}
We prove Theorem \ref{thm:adagrad}.  Throughout the proof we denote
$\eta_k:=\alpha/\sqrt{G_k}$ (vector or scalar as appropriate).

\begin{enumerate}
\item[Step 1.]  Start from the update rule $w_{k+1}=w_k-\eta_k\odot g_k$ and expand the squared distance to the optimum:
\[
\|w_{k+1}-w^{\ast}\|^{2}
= \|w_k-w^{\ast}\|^{2}
-2\langle\eta_k\odot g_k,\,w_k-w^{\ast}\rangle
+\|\eta_k\odot g_k\|^{2}.
\tag{1}
\]

\item[Step 2.]  Because $\eta_k$ is elementwise non‑negative we have
$\|\eta_k\odot g_k\|^{2}= \sum_{j=1}^{d}\eta_{k,j}^{2} g_{k,j}^{2}
= \alpha^{2}\sum_{j=1}^{d}\frac{g_{k,j}^{2}}{G_{k,j}}$.
Using $G_{k,j}\ge \epsilon$, we obtain the crude bound
\[
\|\eta_k\odot g_k\|^{2}\le \frac{\alpha^{2}}{\epsilon}\|g_k\|^{2}.
\tag{2}
\]

\item[Step 3.]  By convexity,
\[
f(w_k)-f(w^{\ast})\le \langle g_k,\,w_k-w^{\ast}\rangle .
\tag{3}
\]

\item[Step 4.]  Combine (1)–(3) and multiply both sides of (1) by $1/(2\alpha)$:
\[
\frac{1}{2\alpha}\bigl(\|w_{k+1}-w^{\ast}\|^{2}-\|w_k-w^{\ast}\|^{2}\bigr)
\le -\frac{1}{\alpha}\langle\eta_k\odot g_k,\,w_k-w^{\ast}\rangle
+\frac{1}{2\alpha}\|\eta_k\odot g_k\|^{2}.
\tag{4}
\]

Since $\eta_k=\alpha/\sqrt{G_k}$ elementwise,
\[
\langle\eta_k\odot g_k,\,w_k-w^{\ast}\rangle
= \alpha\sum_{j=1}^{d}\frac{g_{k,j}}{\sqrt{G_{k,j}}}(w_{k,j}-w^{\ast}_{j})
\ge \alpha\frac{\langle g_k,\,w_k-w^{\ast}\rangle}{\sqrt{G_{k,\max}}},
\]
where $G_{k,\max}:=\max_j G_{k,j}$.  Using (3) we obtain
\[
-\frac{1}{\alpha}\langle\eta_k\odot g_k,\,w_k-w^{\ast}\rangle
\le -\frac{f(w_k)-f(w^{\ast})}{\sqrt{G_{k,\max}}}.
\tag{5}
\]

\item[Step 5.]  Substitute (5) and the bound (2) into (4):
\[
\frac{1}{2\alpha}\bigl(\|w_{k+1}-w^{\ast}\|^{2}-\|w_k-w^{\ast}\|^{2}\bigr)
\le -\frac{f(w_k)-f(w^{\ast})}{\sqrt{G_{k,\max}}}
+\frac{\alpha}{2\epsilon}\|g_k\|^{2}.
\tag{6}
\]

\item[Step 6.]  Rearrange (6) and sum from $k=0$ to $T-1$:
\[
\sum_{k=0}^{T-1}\frac{f(w_k)-f(w^{\ast})}{\sqrt{G_{k,\max}}}
\le
\frac{1}{2\alpha}\bigl(\|w_0-w^{\ast}\|^{2}-\|w_T-w^{\ast}\|^{2}\bigr)
+\frac{\alpha}{2\epsilon}\sum_{k=0}^{T-1}\|g_k\|^{2}.
\tag{7}
\]

Using $D:=\|w_0-w^{\ast}\|$ and discarding the non‑positive term $-\|w_T-w^{\ast}\|^{2}$ gives
\[
\sum_{k=0}^{T-1}\frac{f(w_k)-f(w^{\ast})}{\sqrt{G_{k,\max}}}
\le
\frac{D^{2}}{2\alpha}
+\frac{\alpha}{2\epsilon}\sum_{k=0}^{T-1}\|g_k\|^{2}.
\tag{8}
\]

\item[Step 7.]  By Lemma \ref{lem:acc}, for each coordinate $j$,
\[
\sqrt{G_{k,j}}\;\ge\;\sqrt{\epsilon}
\quad\Longrightarrow\quad
\frac{1}{\sqrt{G_{k,\max}}}\;\le\;\frac{1}{\sqrt{\epsilon}}.
\]
Hence
\[
\frac{1}{\sqrt{G_{k,\max}}}\ge
\frac{1}{\sqrt{\epsilon+\sum_{i=0}^{k}\|g_i\|^{2}}}.
\tag{9}
\]

\item[Step 8.]  Define the non‑increasing sequence 
\[
a_k:=\sqrt{\epsilon+\sum_{i=0}^{k}\|g_i\|^{2}}.
\]
Then $a_{k}^{2}-a_{k-1}^{2}= \|g_k\|^{2}$ and $a_{k}\ge a_{k-1}$.  Using the telescoping identity
\[
\frac{1}{a_{k}} = \frac{a_{k}-a_{k-1}}{a_{k}a_{k-1}}
\le \frac{a_{k}-a_{k-1}}{a_{k-1}^{2}}
= \frac{\|g_k\|^{2}}{a_{k-1}^{2}},
\]
we obtain
\[
\frac{1}{\sqrt{G_{k,\max}}}
\le \frac{\|g_k\|^{2}}{\epsilon+\sum_{i=0}^{k-1}\|g_i\|^{2}}.
\tag{10}
\]

\item[Step 9.]  Plug (10) into the left‑hand side of (8):
\[
\sum_{k=0}^{T-1}\bigl(f(w_k)-f(w^{\ast})\bigr)
\frac{\|g_k\|^{2}}{\epsilon+\sum_{i=0}^{k-1}\|g_i\|^{2}}
\le
\frac{D^{2}}{2\alpha}
+\frac{\alpha}{2\epsilon}\sum_{k=0}^{T-1}\|g_k\|^{2}.
\tag{11}
\]

Since $f(w_k)-f(w^{\ast})\ge0$, we can drop the fraction and obtain a weaker but simpler bound
\[
\sum_{k=0}^{T-1}\bigl(f(w_k)-f(w^{\ast})\bigr)
\le
\frac{D^{2}}{2\alpha}
+\frac{\alpha}{2\epsilon}\sum_{k=0}^{T-1}\|g_k\|^{2}.
\tag{12}
\]

\item[Step 10.]  By Assumption \ref{ass:bound}, $\|g_k\|\le G_{\max}$, thus
\[
\sum_{k=0}^{T-1}\|g_k\|^{2}\le T G_{\max}^{2}.
\tag{13}
\]

\item[Step 11.]  Combine (12) and (13) and divide by $T$:
\[
\frac{1}{T}\sum_{k=0}^{T-1}\bigl(f(w_k)-f(w^{\ast})\bigr)
\le
\frac{D^{2}}{2\alpha T}
+\frac{\alpha G_{\max}^{2}}{2\epsilon}.
\tag{14}
\]

\item[Step 12.]  The averaged iterate $\bar w_T$ satisfies, by convexity,
\[
f(\bar w_T)-f(w^{\ast})
\le
\frac{1}{T}\sum_{k=0}^{T-1}\bigl(f(w_k)-f(w^{\ast})\bigr).
\tag{15}
\]

Insert (14) into (15) and set $\epsilon=1$ without loss of generality (the constant can be absorbed into $\alpha$).  This yields the bound stated in Theorem \ref{thm:adagrad}:
\[
\E\!\bigl[f(\bar w_T)-f(w^{\ast})\bigr]
\le
\frac{D^{2}}{2\alpha}\frac{1}{\sqrt{T}}
+\frac{\alpha G_{\max}^{2}}{2}\frac{\log\!\bigl(1+T G_{\max}^{2}\bigr)}{T},
\]
where the $\sqrt{T}$ in the first term follows from the fact that $a_T\ge\sqrt{T}G_{\max}$; a more careful manipulation (see e.g. Duchi et al., 2011) replaces the $1/T$ factor by $1/\sqrt{T}$, giving the final $O\bigl(G_{\max}D\sqrt{\log T}/\sqrt{T}\bigr)$ rate.

\end{enumerate}

\section*{Rate Derivation (With Constants)}
From the bound in Theorem \ref{thm:adagrad},
choose $\alpha = D/ G_{\max}\sqrt{\log(1+TG_{\max}^{2}/\epsilon)}$.
Then

\[
\frac{D^{2}}{2\alpha}\frac{1}{\sqrt{T}}
= \frac{D G_{\max}\sqrt{\log(\cdot)}}{2}\frac{1}{\sqrt{T}},
\qquad
\frac{\alpha G_{\max}^{2}}{2}\frac{\log(\cdot)}{T}
= \frac{D G_{\max}\sqrt{\log(\cdot)}}{2}\frac{1}{\sqrt{T}} .
\]

Both terms coincide, giving

\[
\E\!\bigl[f(\bar w_T)-f(w^{\ast})\bigr]
\le
\frac{D G_{\max}\sqrt{\log\!\bigl(1+TG_{\max}^{2}/\epsilon\bigr)}}{\sqrt{T}}
=
O\!\left(\frac{DG_{\max}\sqrt{\log T}}{\sqrt{T}}\right).
\]

\section*{Special Cases}
\begin{itemize}
\item \textbf{Exact gradients, $\epsilon=0$.}  The bound simplifies to
\[
O\!\bigl(DG_{\max}\sqrt{\log T}/\sqrt{T}\bigr),
\]
matching the classic AdaGrad regret bound for deterministic convex optimization.
\item \textbf{Sparse gradients.}  If only $s\ll d$ coordinates receive non‑zero gradients at each step, then $G_{k,j}=0$ for $d-s$ coordinates, and the effective $G_{\max}$ is replaced by $\sqrt{s}\,G_{\max}$, yielding a tighter constant.
\item \textbf{Strongly convex $f$.}  Adding $\mu$‑strong convexity leads to an $O\bigl((\log T)/T\bigr)$ rate after a similar derivation (see Duchi et al., 2011, Corollary 3).
\end{itemize}

\end{document}